\section{Introduction}
\label{sec:introduction}

Modern cloud platforms (e.g., AWS IoT Core, Azure IoT Edge) centrally manage edge gateways where router firmware mediates operational telemetry. In cloud DevSecOps pipelines, operators must vet firmware updates for vulnerabilities before deployment. However, edge firmware frequently invokes privileged shell commands via \texttt{system}, \texttt{popen}, and vendor wrappers. When SecOps teams run static taint scanners~\cite{firmalice2015,karonte2020,satc2021,hermescan2024,mango2024,ye2024cmd}, they face hundreds of unranked warnings that cause severe triage fatigue and stall compliance. While rehosting and symbolic execution provide deeper modeling~\cite{firmadyne2016,firmae2020,wright2021rehosting,angr2016,surveyse2018}, full-firmware emulation fails on missing peripherals and state explosion.

The key challenge in command-injection triage is not whether attacker-controlled data reaches a shell sink, but whether attacker-controlled bytes occupy executable shell syntax positions in the final command string. A static closure merely relates a source to a sink without identifying dependent final bytes or proving syntactic activation. This gap stems from three practical challenges:
\emph{(1)~Reachability vs.\ Command-Byte Control:} taint tracking flags dataflow paths but cannot prove that untrusted bytes govern execution tokens.
\emph{(2)~Contextual Quoting:} metacharacters execute only if unescaped by surrounding quote states (unquoted, single, or double quotes).
\emph{(3)~DevSecOps Triage Fatigue:} operators lack verifiable proof receipts, while unreached paths are wrongly treated as safe negatives.

To overcome this bottleneck, we present \tool, an \emph{evidence-driven triage framework} for cloud-managed edge devices. Complementary to upstream static analyzers (e.g., SaTC, Mango, CodeQL), \tool serves as a downstream verification funnel that ingests candidate warnings and evaluates observable syntactic control. Unlike standard targeted symbolic execution that only checks sink reachability, \tool executes a 5-stage triage pipeline: \emph{warning ingestion $\to$ argument binding $\to$ final-byte reconstruction $\to$ quote-aware matrix evaluation $\to$ cryptographic ledger emission}. It enforces four non-escalation admission rules: front-end closure is not final-byte dependence; template feasibility is not source linkage; projected SAT is rechecked under the full path store; and incomplete search yields an explicit typed obligation. \tool evaluates an 11-dimensional shell effect matrix that defaults to fail-closed for unmodeled syntax (Fig.~\ref{fig:contract}).

\paragraph{Two-Tiered Evaluation \& Practical Impact}
We evaluate \tool on two tiers. First, on a 24-case controlled benchmark suite modeling sanitization, quoting, and truncation primitives (inspired by NIST SARD CWE-78), \tool achieves perfect classification accuracy in distinguishing evidence-supported from unsupported claims; this benchmark evaluates evidence classification rather than end-to-end vulnerability discovery, reducing unsupported reachability and regex warnings by isolating exact byte extents and quote contexts. Second, across eight router firmware images spanning \TSDSReviewCampaignRecords{} static warnings, \tool validates \TSDSReviewCampaignCompleteDirect{} command-injection candidates with complete final-byte evidence (alongside \TSDSReviewCampaignObservedPrefixDirect{} observed-prefix candidates), reducing the number of candidates requiring manual semantic inspection by 92.1\% on the evaluated corpus while preserving all \TSDSReviewCampaignResidual{} unresolved cases as explicit typed obligations. Dynamic QEMU emulation confirms that synthesized witnesses reach unquoted sinks to reproduce command execution behaviors corresponding to 6 known CVE patterns and 5 boundary controls, without claiming automated exploit generation. With a median triage time of 9.09~s per warning, \tool establishes an auditable evidence pipeline without rehosting bottlenecks.

\paragraph{Contributions}
We make three contributions to cloud and edge security analysis:
\begin{enumerate}
  \item An executable \emph{final-command-byte evidence model} with explicit admission rules, distinguishing final-byte syntactic control from reachability closures.
  \item \emph{Evidence-directed symbolic execution} combining sink corridors, semantic argument binding, byte provenance, and a quote-aware shell matrix.
  \item An extensive \emph{empirical firmware evaluation} across eight commercial router images (518 warnings), validating 36 candidate profiles and reproducing 6 CVE patterns within an auditable DevSecOps ledger.
\end{enumerate}
